\documentclass[10pt,a4paper]{amsart}
\usepackage[margin=19mm]{geometry}
\usepackage{amsmath,amssymb,mathtools}
\usepackage[scr=rsfs,cal=euler]{mathalfa}
\usepackage{xfrac}
\usepackage[unicode,hidelinks,bookmarks=false]{hyperref}
\newcommand{\ie}{i.e.\ }
\newcommand{\cf}{cf.\ }
\newcommand{\resp}{respectively\ }
\newcommand{\Subsection}{\subsection}
\providecommand{\nobreakdash}{\nobreak\hbox{-}}
\setlength{\emergencystretch}{2em}
\multlinegap0pt
\usepackage{bm}
\usepackage{bbm}
\usepackage{dsfont}
\usepackage{xspace}
\usepackage{cancel}
\usepackage{mathtools}
\usepackage{amsthm}
\makeatletter
\@ifundefined{lem}{\newtheorem{lem}{Lemma}}{}
\makeatother
\newcommand{\A}{\mathcal A}
\newcommand{\al}{\alpha}
\newcommand{\dl}{\delta} 
\newcommand{\ep}{\varepsilon}
\newcommand{\ey}{\frac{1}{2}} 
\newcommand{\gm}{\gamma}
\newcommand{\rr}{\mathbb R}
\newcommand{\tht}{\theta}
\title[Positive energy solutions in the anisotropic Kepler problem ]{Positive energy solutions in the anisotropic Kepler problem with homogeneous potential}
\author{Guowei Yu}
\date{Source: arXiv:2507.10031v2 (CC BY 4.0)}
\begin{document}
\maketitle

\noindent\emph{Source and licence.} Positive energy solutions in the anisotropic Kepler problem with homogeneous potential. Version arXiv:2507.10031v2 (CC BY 4.0), distributed under the Creative Commons Attribution 4.0 International licence, as recorded in the arXiv licence metadata for this exact version and shown on the abstract page https://arxiv.org/abs/2507.10031v2. Full rights notice: licenses/arxiv-2507.10031-CC-BY-4.0.txt.\par\smallskip

\noindent\textit{Editorial scope.} Counted unit: the Lem whose source label is \texttt{lem;rho-bound} in the article (source0.tex lines 736--738, source label \texttt{lem;rho-bound}), delivered together with the article's own complete original proof of it (source0.tex lines 740--824, one \texttt{proof} environment of 85 lines). No mathematics is written, repaired or re-derived: statement and proof are the article's own text line for line. Required context retained uncounted: lem;I-ge (lines 363--368); lem;|s-dot| (lines 392--397). Every definition, display and label of those lines is retained unchanged. Omitted from this package and not redistributed: 1--362, 369--391, 398--735, 739--739, 825--1341. Nothing inside a retained excerpt is omitted or altered: every retained body is delivered exactly as the article's lines are written, so no omission receipt of any kind accompanies this package. Citations: the retained excerpts carry no citation command at all, so no bibliography entry of the article is transcribed and no reference is redistributed. Reference adaptation: none was needed and none was made; every reference inside the delivered unit resolves inside it, so no unresolved reference or citation is produced. Printed slips kept literal: no printed word, symbol, hypothesis or number of the counted statement or of its proof was altered. Layout and class: the article's own class is \texttt{amsart} with packages including mathrsfs, graphicx, latexsym, tikz, color, euscript, amsfonts, amssymb, amsmath, amscd, amsthm, epstopdf, hyperref, upgreek; the wrapper is the lane's pinned \texttt{amsart} editorial wrapper at 10pt on A4, which restates the theorem-like environments the retained text needs and adds the article's own macro definitions that the retained text needs (\texttt{\textbackslash A}, \texttt{\textbackslash al}, \texttt{\textbackslash dl}, \texttt{\textbackslash ep}, \texttt{\textbackslash ey}, \texttt{\textbackslash gm}, \texttt{\textbackslash rr}, \texttt{\textbackslash tht}), with extra wrapper packages (bm, bbm, dsfont, xspace, cancel, mathtools). The delivered numbering is the unit's own. Assets: no retained span contains a figure environment, a graphics-inclusion command, a PDF-page inclusion, a table or a TikZ picture; no image file is copied into this package, the package declares no asset dependency and no image file is redistributed with it. Licence: version arXiv:2507.10031v2 of this article is distributed under the Creative Commons Attribution 4.0 International licence, as recorded in the arXiv licence metadata for this exact version and shown on the abstract page https://arxiv.org/abs/2507.10031v2. A full rights notice is delivered at licenses/arxiv-2507.10031-CC-BY-4.0.txt. Characters: every retained excerpt is pure ASCII, so the delivered engine, XeLaTeX with its default Latin Modern text font, reports no missing character.
\begin{lem}
\label{lem;I-ge} For any $T^-<t_0< t_1 <T^+$, 
\begin{equation}
\label{eq;I-ge} I(t_1) \ge 2h(t_1 -t_0)^2 + \dot{I}(t_0)(t_1 -t_0) + I(t_0).  
\end{equation}
\end{lem}

\begin{lem}
\label{lem;|s-dot|} If $\dot{I}(t_0) \ge 0$ for some $t_0 \in (T^-, T^+)$, then for any $t_0 < t_1 < t_2 < T^+$, 
\begin{equation}
\label{eq;int-s-dot} \int_{t_1}^{t_2} |\dot{s}| \,dt \le  \frac{2}{\al} \frac{\sqrt{2h r^{\al}(t_0)+2U_{max}}}{(2h )^{\frac{2 +\al}{4}}}(t_1-t_0)^{-\frac{\al}{2}} .
\end{equation}
\end{lem}

\begin{lem} \label{lem;rho-bound}
When $|\tht_+ -\tht_-| >\pi$ and each $\gm_n$ is collision-free, $\liminf_{n \to \infty} \rho_n <\infty.$
\end{lem}

\begin{proof}
By a contradiction argument, after passing to a  sub-sequence, we may assume 
$$ \lim_{n \to \infty} \rho_n =\infty \text{ and } d = \lim_{n \to \infty} \rho_n /R_n \in [0, 1]. $$

%\begin{equation}
%\label{eq;I-n-t} |\gm_n(t)|^2 \ge 2h |t|^2 + \rho_n^2, \; \forall t \in [T_n^-, T_n^+].
%\end{equation}
First let's consider the case that $d >0$. By Lemma \ref{lem;I-ge}, $R_n^2 \ge 2h (T^{\pm}_n)^2 + \rho_n^2$. This implies
$$ |T_n^{\pm}| \le \frac{1}{\sqrt{2h}} \left( 1 -\left( \frac{\rho_n}{R_n} \right)^2 \right)^{\ey} R_n. $$
By the energy identity,
\begin{equation*} 
\begin{aligned}
\A_h(\gm_n; T_n^-, T_n^+) & = 2 \int_{T_n^-}^{T_n^+} h + U(\gm_n(t)) \,dt \le 2 \left( h + \frac{U_{max}}{\rho_n^{\al}} \right) (T^+_n -T^-_n) \\
& \le \frac{4}{\sqrt{2h}} \left( h + \frac{U_{max}}{\rho_n^{\al}} \right)  \left( 1 -\left( \frac{\rho_n}{R_n} \right)^2 \right)^{\ey} R_n
\end{aligned}
\end{equation*} 
Our assumptions then imply the existence of a $\dl \in (0, 1)$, such that when $n$ is large enough,
\begin{equation} \label{eq;Ah-xn>Rn}
\A_h(\gm_n; T_n^-, T_n^+)  \le 2 \sqrt{2h} (1 - \dl) R_n.
\end{equation}
Meanwhile as $ |\tht_+ - \tht_-|>\pi$, we can find two instants $T_n^- < t_n^- \le t_n^+ < T_n^+$ satisfying
$$ \langle \gm_n(t_n^{\pm}), \gm_n(T_n^{\pm}) \rangle =0,$$
Then the triangle inequality implies 
$$ |\gm_n(T_n^{\pm}) - \gm_n(t_n^{\pm})| \ge |\gm_n(T_n^{\pm}) -0| \ge R_n. $$
Using this, we can get 
$$  
\A_h(\gm_n; T_n^-, T_n^+)  \ge \int_{T_n^-}^{t_n^-} + \int_{t_n^+}^{T_n^+} \ey |\dot{\gm}_n|^2 + h \,dt  \ge \sqrt{2h}   \left( \int_{T_n^-}^{t_n^-}  + \int_{t_n^+}^{T_n^+} |\dot{\gm}_n|  \,dt \right) \ge 2 \sqrt{2h} R_n.  
$$
This is a contradiction to \eqref{eq;Ah-xn>Rn}. 

Now let's consider the case $d=0$. For each $n$, set 
$$ y_n(\tau) = \rho_n^{-1} \gm_n(\rho_n \tau), \; \forall \tau \in [T_n^-/\rho_n, T_n^+/\rho_n]. $$
By a direct computation,
\begin{equation}
\label{eq;yn''} \begin{cases}
& y''_n(\tau) = \frac{\nabla U(y_n(\tau))}{\rho_n^{\al}  }; \\
& \ey |y'_n(\tau)|^2 - \frac{U(y_n(\tau))}{\rho_n^{\al} } = h.
\end{cases}
\end{equation} 

After passing to a proper sub-sequence, $y_n \to y_{\infty}$ in $C^2_{\text{loc}}(\rr, \rr^2)$, where 
$$ \begin{cases}
y''_{\infty}(\tau) & = 0; \\
\ey |y_{\infty}'(\tau)|^2 & = h. 
\end{cases}
$$
Since $|y_n(0)| = 1$, we may further assume $y_n(0) \to e^{i\tht_0}$, as $n \to \infty$, for some $\tht_0 \in \rr$. Then  
\begin{equation*}
\label{eq;y-infty}  y_{\infty}(\tau) =  \begin{cases}
e^{i \tht_0} + \sqrt{2h} e^{i(\tht_0 +\frac{\pi}{2})} \tau, & \text{ if } \tht_- < \tht_+; \\
 e^{i \tht_0} + \sqrt{2h} e^{i(\tht_0 - \frac{\pi}{2})} \tau, & \text{ if } \tht_- > \tht_+.  
\end{cases}   
\end{equation*}
As a result, 
\begin{equation}
\label{eq;s-infty} \lim_{\tau \to + \infty} |\text{Arg}(y_{\infty}(\tau)) - \text{Arg}(y_{\infty}(-\tau))| = \pi. 
\end{equation}

To get a contradiction, set 
$$ \tht_{\infty}(\tau)= \text{Arg}(y_{\infty}(\tau)), \; \tht_n(\tau)= \text{Arg}(y_n(\tau)). $$ 
Then for any $\ep>0$ and $\tau_1$, there is an integer $n_0= n_0(\ep, \tau_1)$, such that   
\begin{equation*}
|\tht_n(\tau_1)-\tht_{\infty}(\tau_1)| \le \ep/2, \; \forall n \ge n_0. 
\end{equation*}
Meanwhile $\lim_{n \to \infty} \rho_n = \infty$ implies the existence of a $\tau_{\ep}>0$ large enough, such that 
\begin{equation*}
\label{eq;tau-ep} \frac{2}{\al} \frac{\sqrt{2h + 2 U_{max} \rho_n^{-\al}}}{(2h)^{\frac{2+\al}{4}}} \tau_{\ep}^{-\frac{\al}{2}} \le \ep/2, \;\; \forall n. 
\end{equation*}

Let  $s_n(t) = \gm_n(t)/|\gm_n(t)|$. By a direct computation, 
$$ \int_{\tau_{\ep}}^{T_n^+/\rho_n} |\tht'_n(\tau)| \,d\tau = \int_{\rho_n \tau_{\ep}}^{T^+} |\dot{s}_n(t)| \,dt 
. $$
Then Lemma \ref{lem;|s-dot|} implies that for any $\tau_1 \ge \tau_{\ep}$,
$$ |\tht_n(\tau_1)-\tht_+| \le \int_{\tau_{\ep}}^{T_n^+/\rho_n} |\tht'_n(\tau)| \,d\tau \le \frac{2}{\al} \frac{\sqrt{2h + 2 U_{max} \rho_n^{-\al}}}{(2h)^{\frac{2+\al}{4}}} \tau_{\ep}^{-\frac{\al}{2}} \le \ep/2.
$$

The above estimates show
$$ |\tht_{\infty}(\tau_1) - \tht_+| \le |\tht_{\infty}(\tau_1) -\tht_n(\tau_1)| + |\tht_n(\tau_1)- \tht_+| \le \ep. $$
As this holds for any $\ep>0$ small and $\tau_1 \ge \tau_{\ep}$, we get  
$$ \lim_{\tau_1 \to +\infty} \tht_{\infty}(\tau_1)=\tht_+. $$
A similar argument as above shows 
$$ \lim_{\tau_1 \to -\infty} \tht_{\infty}(\tau_1)=\tht_-. $$
As a result, $\lim_{\tau \to +\infty}|\tht_{\infty}(\tau) - \tht_{\infty}(-\tau)| = |\tht_+ -\tht_-| >\pi$, which is a contradiction to \eqref{eq;s-infty}.

\end{proof}

\end{document}
